\itemhead{习题4.2 · 第4题（教材p.113）}{XJ-B1-C04-HW-P113-Y1315}
\label{XJ-B1-C04-HW-P113-Y1315}
{\footnotesize 湘教版教材p.113 / PDF第120页 / 习题4.2 / 第4题}
\field{原题}{食品的保鲜时间 $y$（h）与温度 $x$（$^\circ\mathrm C$）的关系为 $y=\mathrm e^{kx+b}$，其中 $\mathrm e=2.71828\cdots$，$k,b$ 为常数。某种食品在 $0\,{}^\circ\mathrm C$ 时保鲜时间为192 h，在 $22\,{}^\circ\mathrm C$ 时保鲜时间为48 h，求该食品在 $33\,{}^\circ\mathrm C$ 时的保鲜时间。}
\field{参考答案}{由 $e^b=192,e^{22k}=48/192=1/4$，得 $e^{33k}=(1/4)^{3/2}=1/8$；故33摄氏度时为24小时。无需使用对数。}
{\small \field{关联知识}{E03 指数增长、衰减与实际模型（教材帮p.131、132、154 / PDF140、141、163）；P02 有理数、实数指数幂与运算（教材帮p.144、145 / PDF153、154）}
\field{方法与易错点}{确认初始量、时间单位、倍率与时刻起点，再代入或逐次比较阈值。 易错：把年折旧10\%写成乘0.1；“超过”用成“达到”；时间起点错一位；把衰减至接近0当实际为0。}
\field{答案依据}{答案册相关题组：PDF第32页（册内页31）；按当前题库原题逐项重算，不按旧版题号机械对套。}}
